\chapter{Introducción}
\label{chap:introduccion}

\lettrine{E}{n} este primer capítulo se presentan los datos básicos de la práctica, una descripción general de su temática y los objetivos principales y secundarios de la misma.

\section{Datos de la práctica}
\label{sec:datospractica}
\begin{itemize}
    \item {\bf Empresa:} CINFO
    \item {\bf Duración:} 23 Septiembre 2024 - 2 Diciembre 2024
    \item {\bf Lugar de realización:} A Coruña
    \item {\bf Tipo:} Modalidad presencial
    \item {\bf Departamento/equipo:} Equipo de IA
    \item {\bf Cargo responsable de práctica:} IA Developer Leader
\end{itemize}

\section{Objetivos}
\label{sec:objetivos}

Las prácticas se centraron en la mejora de sistemas de visión artificial y aprendizaje automático para el reconocimiento de marcadores deportivos en tiempo real, incluyendo la configuración y optimización de herramientas existentes, así como en la investigación e implementación de sistemas de síntesis de voz (TTS).

\subsection{Objetivos generales}
\begin{enumerate}
    \item Familiarización con el flujo de trabajo de machine learning de la empresa
    \item Mejorar el sistema de detección y reconocimiento de marcadores deportivos
    \item Investigar y desarrollar sistemas de síntesis de voz para narración deportiva
    \item Contribuir al desarrollo de herramientas de automatización deportiva
\end{enumerate}

\subsection{Objetivos específicos}
\begin{enumerate}
    \item Optimizar la detección de marcadores digitales en diferentes condiciones
    \item Implementar mejoras en la interfaz de configuración de marcadores
    \item Evaluar y comparar diferentes modelos de síntesis de voz
    \item Crear datasets y realizar fine-tuning de modelos TTS
    \item Ejecutar y evaluar experimentos utilizando MLflow
    \item Documentar y mantener el código desarrollado
\end{enumerate}